\documentclass[11pt]{article}
\usepackage{amsmath,amssymb}
\usepackage[hidelinks]{hyperref}
\title{An optical rendezvous with two outstanding reports}
\date{24 November 2026}
\begin{document}
\maketitle
\paragraph{Plant and local histories.}
Use a five-node indoor inspection model: base \(0\), hub \(1\),
two corridor nodes \(2,3\), and refuge \(4\). The links are
\(01,12,13,24,34\). Exactly one last link, \(24\) or \(34\), is
blocked. The four unblocked possible edges are available for one-tick
moves, but opposite traversals of one edge on the same tick cannot
both take effect. Both crews can occupy a node. An earlier
communication failure has left no radio channel. The two crews
begin with distinct existing duties, not merely locations to visit.
Crew A at \(4\) holds a report, may travel at most two edges, and
must end back at \(4\). Crew B at \(1\) must deliver its own report
to base by tick one, later take and deliver A's report, and end
at base by the final deadline. An actual report handoff takes a
full tick stationary at a common node, with compatible give and take
actions. Co-location does not by itself deliver a report.

An actor observes its own node, its local incident links, its own
custody and received values, and a co-located peer. In particular
A knows the blocked link at its refuge; B does not know it at the hub.
A and B may choose an optical mode at nodes \(4\) and \(1\),
respectively. One simultaneous tick of beacon and listening
produces a source-coded branch observation at B. The actuator and
receiver actions use their tick without travelling. This is a stipulated
one-way optical pathway, not inferred from the movement links.
The possible-erasure version includes both worlds where this observation
can fail on every attempt. No surprise resource or other channel is
made available. The plant is a finite development surrogate; the
optical geometry, clocks, custody sensing, and actual contact duration
would require separate source validation.

\paragraph{Exact deadline.}
B must move from hub to base during the first tick, so the earliest
return to the hub is tick two. With the branch already known, it can
reach the open corridor node at tick three, transfer through tick four,
and return to base by tick six. A can reach that node and return to
refuge with its two travel units. No shorter delivery is possible:
B must pass through the hub, reach a meeting node, pay the dwell,
and return to base. A cannot travel as far as the hub and still return
within budget.

With no source-coded signal, departure through either candidate branch
at tick two is chosen from B's identical views. The adverse branch
takes B to the wrong corridor node at tick three. B cannot meet A there.
Its shortest continuation goes back through hub at tick four, to the
other node at tick five, then transfers and returns to base by tick
eight. Waiting at the hub instead cannot expose the blockage, and
A cannot reach that hub and return. Taking the first branch and
revising on the observed closed link attains the same bound.

With reliable optics, B can listen at the hub during the tick beginning
at two. The new observation is available at tick three. This enables
choice of the correct branch and completion at tick seven. To finish
by six B would have to leave the hub at tick two already knowing which
branch is open. It cannot have listened during the first tick because
its old report is due at one, and cannot listen at base. Thus seven is
tight. The construction can encode the bit by a pre-agreed pulse or
certified silence; silence loses that meaning when an emitted pulse
can be erased. If erasure is allowed in both worlds, the all-erased
histories give the eight-tick lower bound. Ignoring the optical
mechanism attains eight in all histories.

\paragraph{Comparison.}
The direct exact search assigns one action per actor per local-history
class before applying it to the source worlds. Partition refinement
only adds the observations actually received and never forgets a prior
distinction. A separate rollout executes the extracted tables using
only each actor's local observation history and prior moves.
The written deadline proof and the exact search agree across 30
deadline checks, and successful optical and erased-channel policies
replay. The matched planner is an exact solution of this finite model;
there is no claim that the optical policy beats it.
The result is a test of a limited source substitution, with an actual
old duty that prevents using the direct channel one tick earlier.
It is not a held-out test, broad transfer, new historical result, or
acceptance of R1, R3a, or I2.

\paragraph{Frozen receiving controller under a physical supplier change.}
The pulse requires a sending action at the reading cut itself.
A second source instead permits A to set a mechanical flag at the
refuge. A's earlier work prohibits issuing this command before tick
\(u\). Issuance occupies one tick, and the flag is physically
available only \(\ell\geq 1\) ticks after issuance. The flag remains
latched until B reads it at the hub. In the model a latch damaged
before the reading produces no bit, even if its command was executed.
The source implementations differ in when the decisive action must
be taken and in which physical event constitutes the usable output.

We froze B's original local-history policy (11 entries) from the
reliable optical case, before running the mechanical source.
On the branch originally represented by silence, A leaves the flag
unset. On the other branch the unchanged policy observes the necessary
bit at the listening cut, tick two, if and only if the flag remains
healthy and \(u+\ell\leq 2\). The earliest permitted issue realizes
that condition, and every schedule starting later is also too late
if that earliest one fails. A's retained report still has to be
physically transferred after the cue; producing it is not delivering
the report. Source replay checks the actual later transfer and
both safe terminal positions, not just report equality. In the
example, issuing at tick zero can accommodate a two-tick lag.
The same flag after one more tick of prior work misses the cut.

The 84 scope-preservation comparisons use the byte-identical
receiver table, whose digest is pinned in the checker. Successful
adaptation requires no receiver rewrite. The timing inequality is
also the equally factored direct comparison; no advantage in proof
work is measured or claimed. This was an authored source selection
after freezing a finite policy, not an independently selected
held-out replacement. The latch issuing discipline, physical
contact, and accepted readout meanings are model premises.
\end{document}
